\subsection{Цели, задачи и требования к разрабатываемому решению}

\subsubsection{Постановка цели исследования}

Анализ предметной области, проведённый в разделах \ref{subsec:relevance-research} и \ref{subsec:review-existing-solutions},
позволяет сформулировать основное противоречие, определяющее
актуальность настоящей работы. С одной стороны, механизм work-stealing
обладает строгим теоретическим обоснованием и демонстрирует высокую
эффективность в промышленных реализациях. С другой стороны,
существующие промышленные реализации (Intel TBB, Microsoft PPL)
представляют собой крупные библиотеки со сложным внутренним
устройством, не допускающим лёгкой модификации и систематического
исследования. Стандартная библиотека C++ (вплоть до C++17) не
предоставляет готового пула потоков с механизмом work-stealing,
а учебные реализации не подкреплены экспериментальными данными
о производительности.

В связи с изложенным целью настоящей выпускной квалификационной
работы является разработка библиотеки на языке C++, реализующей
пул потоков с механизмом work-stealing, и экспериментальное
исследование её производительности в сравнении с классическим
пулом потоков при различных сценариях нагрузки.

Поставленная цель предполагает решение как инженерной, так и
исследовательской задачи. Инженерная составляющая состоит в
проектировании и реализации библиотеки, построенной исключительно
на стандарте языка версии C++17, без внешних
зависимостей. Исследовательская составляющая состоит в проведении
серии измерений производительности, позволяющих количественно
оценить преимущества и ограничения механизма work-stealing при
различных условиях применения.

Выбор стандарта C++17 в качестве технологической основы
обусловлен тремя соображениями. Во-первых, все необходимые примитивы
многопоточности --- \texttt{std::thread}, \texttt{std::atomic},
\texttt{std::mutex}, \texttt{std::condition\_variable},
\texttt{std::future} --- были введены в стандарте C++11 и в полной
мере доступны начиная с этой версии языка. Во-вторых, стандарт
C++17 по-прежнему остаётся основным в большинстве
производственных кодовых баз, что обеспечивает широкую применимость
разрабатываемой библиотеки. В-третьих, основная методическая база
настоящей работы --- книга Уильямса «C++ Concurrency in
Action»~\cite{Williams2019} --- рассматривает многопоточность
именно в контексте данных стандартов, что обеспечивает строгое
соответствие между теоретической основой и практической реализацией.

\subsubsection{Декомпозиция цели на конкретные задачи}

Для достижения поставленной цели необходимо решить следующие задачи.

Первая задача --- анализ существующих подходов и решений.
Провести систематический обзор теоретических работ по механизму
work-stealing, классических моделей управления потоками и существующих
реализаций пулов потоков. Выявить достоинства и недостатки
существующих решений, обосновать требования к разрабатываемой
библиотеке. Данная задача решается в главе~1 настоящей работы.

Вторая задача --- проектирование архитектуры библиотеки.
Разработать архитектуру библиотеки, включающую две реализации
пула потоков: классический пул потоков с единой глобальной очередью
задач и пул потоков с механизмом work-stealing. Спроектировать
единый программный интерфейс, позволяющий использовать обе
реализации взаимозаменяемо и проводить их сравнительное
исследование в одинаковых условиях.

Третья задача --- реализация потокобезопасной двусторонней очереди. 
Разработать и реализовать структуру данных (deque) для механизма work-stealing. 
Реализация должна обеспечивать корректную синхронизацию доступа владельца очереди и потоков-«воров» при краже задач, гарантируя отсутствие гонок данных посредством блокирующей синхронизации (std::mutex).

Четвёртая задача --- реализация пулов потоков.
На основе спроектированной архитектуры реализовать классический
пул потоков с глобальной очередью и пул потоков с механизмом
work-stealing. Обе реализации должны быть построены исключительно
на стандартных средствах C++17 и не иметь внешних
зависимостей.

Пятая задача --- тестирование корректности.
Разработать набор тестов, верифицирующих корректность работы
реализованных структур данных и пулов потоков. Тестирование должно
охватывать граничные случаи: одновременный доступ к очереди
владельца и вора, поведение при пустой очереди, корректность
завершения работы пула при наличии незавершённых задач.

Шестая задача --- экспериментальное исследование производительности.
Разработать методику и набор тестов производительности (benchmarks),
позволяющих сравнить классический пул потоков и пул потоков
с work-stealing при различных сценариях нагрузки. Исследовать
зависимость производительности от числа рабочих потоков, размера
и характера задач. На основе полученных данных сформулировать
выводы об условиях эффективного применения механизма work-stealing.

Седьмая задача --- разработка документации и рекомендаций.
Подготовить документацию к разработанной библиотеке и
сформулировать практические рекомендации по её применению
на основе результатов экспериментального исследования.

В ходе анализа предметной области и обзора существующих решений
были получены следующие выводы.

Во-первых, переход к многоядерным архитектурам сделал параллелизм
неотъемлемым требованием к современному программному обеспечению.
Классические подходы к управлению потоками --- модель «поток на
задачу» и пул потоков с единой глобальной очередью --- обнаруживают
существенные ограничения с точки зрения масштабируемости
и эффективности балансировки нагрузки при большом числе ядер.

Во-вторых, механизм work-stealing обладает строгим теоретическим
обоснованием и доказанной оптимальностью по времени выполнения
и потреблению памяти. Вместе с тем существующие промышленные
реализации (Intel TBB, Microsoft PPL) не допускают лёгкого изучения
и модификации, а стандартная библиотека C++ не предоставляет
готовой реализации пула потоков с work-stealing.

В-третьих, сформулированные цель и задачи работы однозначно
определяют её границы и критерии успеха: разработать библиотеку
на чистом C++17 без внешних зависимостей, реализующую обе
модели пула потоков с единым интерфейсом, и провести их
сравнительное экспериментальное исследование.